\chapter{提案手法}\label{chap:method}

\section{数式}

インライン数式 $E = mc^2$ や $\vect{x} \in \R^n$ が書ける．
別行立て数式：
\begin{equation}
  e^{i\pi} + 1 = 0
  \label{eq:euler}
\end{equation}
複数行の数式：
\begin{align}
  f(x) &= \sum_{k=0}^{\infty} \frac{x^k}{k!} \label{eq:exp} \\
  \diff{f}{x} &= f(x)
\end{align}
番号なし：
\[
  \int_0^1 x^2 \, \mathrm{d}x = \frac{1}{3}
\]
式\eqref{eq:exp}は指数関数である．

\section{図と表}

\begin{figure}[htbp]
  \centering
  \includegraphics[width=0.5\linewidth]{figures/sample.png}
  \caption{サンプル図}
  \label{fig:sample}
\end{figure}

\begin{table}[htbp]
  \centering
  \caption{実験結果}
  \label{tab:result}
  \begin{tabular}{lcc}
    \toprule
    手法 & 精度 & 時間 [s] \\
    \midrule
    従来手法 & 0.85 & 12.3 \\
    提案手法 & \textbf{0.92} & 8.1 \\
    \bottomrule
  \end{tabular}
\end{table}
